% Scratch file for Claude's computations, derivations, and intermediate work.
% Review results here before incorporating into the main writeup.

% ============================================================
% Professional rewrite of Section 7.1, lines 11-23
% (Recap: Perturbative String Theory and Higher-Derivative Supergravity)
% ============================================================

\section{Scratch}
\subsection{Recap: Perturbative String Theory and Higher-Derivative Supergravity}
Higher-derivative corrections to supergravity are most efficiently obtained not from Feynman diagrams, but from the low-energy expansion of string amplitudes.
The perturbative expansion of string amplitudes is organized by two parameters: $\alpha'$, governing the derivative expansion of the low-energy effective action, and $g_s$, weighting the genus expansion of the worldsheet. An $n$-point amplitude admits the expansion
\begin{equation}
    \mathcal A_n = \sum_{g=0}^\infty g_s^{2g-2+n}  \mathcal A_n^{(g)}(\alpha').
\end{equation}

% ============================================================
% Verification of Section 6.2 (Covariant Unification in 12d):
%   (1) sign in the definition of \tilde{\hat F}_7
%   (2) coefficient of the Chern-Simons term in the "12d" action
% Conventions (as used in the section, validated by the |F_4|^2 check):
%   tau_1 = C, tau_2 = e^{-Phi};
%   \hat g_{ab} = (1/tau_2) [[1, tau_1],[tau_1, |tau|^2]], det \hat g_{ab} = 1;
%   \hat C_3 = B_2 ^ dz_1 + C_2 ^ dz_2,  \hat F_4 = H_3 ^ dz_1 + F_3 ^ dz_2;
%   \hat F_7 = F_5 ^ dz_1 ^ dz_2;
%   section's composite: \tilde F_5 = F_5 - (1/2) C_2^H_3 + (1/2) B_2^F_3;
%   IIB CS term (2__IIB.tex): -(1/(4 kappa^2)) \int C_4 ^ H_3 ^ F_3;
%   kappa relation: Vol(T_2)/kappa_{12}^2 = 1/kappa^2.
% ============================================================

\subsection*{Verification: $\tilde{\hat F}_7$ sign and the 12d Chern--Simons coefficient}

\paragraph{Wedge lemma.}
With $\hat C_3 = B_2\wedge dz_1 + C_2\wedge dz_2$ and $\hat F_4 = H_3\wedge dz_1 + F_3\wedge dz_2$, the
terms with a repeated $dz_a$ vanish, leaving
\begin{equation}
\hat C_3\wedge\hat F_4
= B_2\wedge dz_1\wedge F_3\wedge dz_2 + C_2\wedge dz_2\wedge H_3\wedge dz_1 .
\end{equation}
Using $dz_1\wedge F_3 = -F_3\wedge dz_1$ and $dz_2\wedge H_3 = -H_3\wedge dz_2$, then $dz_2\wedge dz_1=-dz_1\wedge dz_2$,
\begin{equation}
B_2\wedge dz_1\wedge F_3\wedge dz_2 = - B_2\wedge F_3\wedge dz_1\wedge dz_2,
\qquad
C_2\wedge dz_2\wedge H_3\wedge dz_1 = + C_2\wedge H_3\wedge dz_1\wedge dz_2,
\end{equation}
so that
\begin{equation}
\boxed{\;\hat C_3\wedge\hat F_4 = \big(C_2\wedge H_3 - B_2\wedge F_3\big)\wedge dz_1\wedge dz_2 .\;}
\label{scratch-wedge}
\end{equation}

\paragraph{(1) Sign of $\tilde{\hat F}_7$.}
With $\hat F_7 = F_5\wedge dz_1\wedge dz_2$ and the section's definition $\tilde{\hat F}_7 = \hat F_7 + \tfrac12\hat C_3\wedge\hat F_4$,
\begin{equation}
\tilde{\hat F}_7
= \Big(F_5 + \tfrac12 C_2\wedge H_3 - \tfrac12 B_2\wedge F_3\Big)\wedge dz_1\wedge dz_2 .
\end{equation}
The bracket is \emph{not} $\tilde F_5 = F_5 - \tfrac12 C_2\wedge H_3 + \tfrac12 B_2\wedge F_3$; the two correction
terms carry the opposite relative sign, so $|\tilde{\hat F}_7|^2 \ne |\tilde F_5|^2$ (the cross term $2F_5\!\cdot\!(\dots)$ flips).
Choosing instead
\begin{equation}
\boxed{\;\tilde{\hat F}_7 \equiv \hat F_7 - \tfrac12\hat C_3\wedge\hat F_4\;}
\end{equation}
gives the bracket $F_5 - \tfrac12 C_2\wedge H_3 + \tfrac12 B_2\wedge F_3 = \tilde F_5$ exactly.
\textbf{Conclusion: the ``$+$'' should be ``$-$''. This is convention-independent (fixed by the section's own $\tilde F_5$).}

\paragraph{(2) Chern--Simons coefficient.}
A 4-form wedged with itself keeps only the cross term ($b\wedge a = (-1)^{16}a\wedge b = a\wedge b$):
\begin{equation}
\hat F_4\wedge\hat F_4 = 2 (H_3\wedge dz_1)\wedge(F_3\wedge dz_2)
= -2 H_3\wedge F_3\wedge dz_1\wedge dz_2 ,
\end{equation}
hence, with $\hat C_4 = C_4$ (so $\hat F_5 = d\hat C_4 = F_5$),
\begin{equation}
\hat C_4\wedge\hat F_4\wedge\hat F_4 = -2 C_4\wedge H_3\wedge F_3\wedge dz_1\wedge dz_2 .
\end{equation}
Integrating ($\int_{T_2}dz_1\wedge dz_2 = \pm Vol(T_2)$, sign set by torus orientation),
\begin{equation}
\tfrac12\int_{T_2\times\mathbb R^{1,9}}\hat C_4\wedge\hat F_4\wedge\hat F_4
= \mp Vol(T_2)\int_{\mathbb R^{1,9}} C_4\wedge H_3\wedge F_3 .
\end{equation}
So the section's display identity is correct \emph{in magnitude} (the $\tfrac12$ cancels the $2$); only its overall
sign is orientation-dependent. The problem is the action coefficient: feeding $\int\hat C_4\hat F_4\hat F_4 = 2 Vol(T_2)\int C_4 H_3 F_3$
into the action term $-\tfrac{1}{4\kappa_{12}^2}\int\hat C_4\hat F_4\hat F_4$ gives
\begin{equation}
-\frac{1}{4\kappa_{12}^2}\cdot 2 Vol(T_2)\int C_4\wedge H_3\wedge F_3
= -\frac{Vol(T_2)}{2\kappa_{12}^2}\int = -\frac{1}{2\kappa^2}\int C_4\wedge H_3\wedge F_3 ,
\end{equation}
which is \emph{twice} the IIB value $-\tfrac{1}{4\kappa^2}\int C_4\wedge H_3\wedge F_3$.
\textbf{Resolution (chosen orientation $dz_2\wedge dz_1 > 0$, i.e. $\int_{T_2}dz_1\wedge dz_2 = -Vol(T_2)$):}
With this orientation $\int\hat C_4\hat F_4\hat F_4 = +2 Vol(T_2)\int C_4\wedge H_3\wedge F_3$, so to reproduce the IIB term
$-\tfrac{1}{4\kappa^2}\int C_4\wedge H_3\wedge F_3$ the action prefactor must be $-\tfrac{1}{8\kappa_{12}^2}$ (replacing $-\tfrac{1}{4\kappa_{12}^2}$):
the factor of $2$ from $\hat F_4\wedge\hat F_4$ halves $\tfrac14\to\tfrac18$, and this orientation keeps the conventional minus sign.
The display identity reads $\tfrac12\int\hat C_4\hat F_4\hat F_4 = +Vol(T_2)\int C_4\wedge H_3\wedge F_3$.
With this orientation the Hodge dual flips sign, $*_{12}F_5 = -(*_{10}F_5)\wedge dz_1\wedge dz_2 = -\hat F_7$ for self-dual $F_5$,
so the self-duality relation reads $\hat F_7 = -*_{12}\hat F_5 \Leftrightarrow F_5 = *_{10}F_5$ (and $*_{12}d\hat C_4 = -\hat F_7$).
All fixes applied to the main text. Independent of any external $|F|^2$ normalization (the $|\hat F_4|^2$ kinetic check matched exactly).

% ============================================================
% Verification of Section 6.3 caveat:
%   the F1/D1/NS5/D5 uplifts do NOT solve the 12d Einstein equation.
%   Clean proof via the Stelle (Sec 6.1) electric-solution warp factors,
%   avoiding a direct 12d Ricci computation.
% ============================================================

\subsection*{Verification: brane uplifts are not 12d solutions (warp-exponent argument)}

\paragraph{Genuine 12d electric brane.}
The 12d theory of Sec.~6.2 is pure gravity plus the 4-form $\hat F_4$, with \emph{no} dilaton.
The Sec.~6.1 electric solution with $D=12$, an $n=4$ form (so worldvolume $d=n-1=3$, transverse
$\tilde d = D-d-2 = 7$) and $\alpha=0$ has
\begin{equation}
\Delta = \alpha^2 + \frac{2 d\tilde d}{D-2} = \frac{2\cdot 3\cdot 7}{10} = \frac{42}{10},
\end{equation}
hence warp exponents
\begin{equation}
-\frac{4\tilde d}{\Delta(D-2)} = -\frac{28}{42} = -\frac23
\quad(\text{worldvolume } t,x_1,z_1),
\qquad
+\frac{4 d}{\Delta(D-2)} = \frac{12}{42} = \frac27
\quad(\text{9 transverse}),
\end{equation}
so a genuine 12d electric 3-brane sourced by $\hat F_4$ would be
\begin{equation}
ds^2 = H^{-2/3}\big(-dt^2+dx_1^2+dz_1^2\big) + H^{2/7}\big(dx_8^2 + dz_2^2\big),
\qquad \nabla^2_{9}H = 0 .
\end{equation}

\paragraph{The F1 uplift differs.}
The F1 uplift of the main text is
\begin{equation}
ds^2 = H^{-3/4}(-dt^2+dx_1^2) + H^{1/4}dx_8^2 + H^{-1/2}dz_1^2 + H^{1/2}dz_2^2 .
\end{equation}
Comparing warps:
\begin{center}
\begin{tabular}{l|cc}
direction & genuine 12d 3-brane & F1 uplift\\\hline
$t,x_1$ & $H^{-2/3}$ & $H^{-3/4}$\\
$z_1$ & $H^{-2/3}$ & $H^{-1/2}$\\
$x_8$ (transv.) & $H^{2/7}$ & $H^{1/4}$\\
$z_2$ & $H^{2/7}$ & $H^{1/2}$
\end{tabular}
\end{center}
A genuine solution warps $(t,x_1,z_1)$ \emph{uniformly} by $H^{-2/3}$; the uplift splits them
($H^{-3/4}$ vs $H^{-1/2}$). Hence the F1 uplift does \emph{not} solve the 12d Einstein equation
$R_{MN}=S_{MN}[\hat F_4]$.

\paragraph{What it is instead.}
The F1 uplift is the 10d F1 (Einstein-frame) solution times a unit-area torus,
$g_{z_1z_1}g_{z_2z_2}=H^{-1/2}H^{1/2}=1$, whose radius ratio $g_{z_1z_1}/g_{z_2z_2}=H^{-1}=e^{2\phi}$
geometrizes the dilaton; the $H^{-3/4}/H^{1/4}$ warps are exactly the 10d F1 warps. So it solves
the \emph{10d} equations of motion, repackaged. The same holds for D1/NS5/D5, which are obtained
from F1 by the torus $S$-rotation and the 12d Hodge duality (the orbit
F1 $\to$ D1 $\to$ NS5 $\to$ D5).
\textbf{Note:} the earlier ``$R_{\mu\nu}=S_{\mu\nu}$ with $n=3,D=10$'' phrasing is exact only for
constant dilaton (flat torus, $R^{(12)}_{\mu\nu}=R^{(10)}_{\mu\nu}$); with the dilaton geometrized
into a warped torus it is not a literal identity, which is why the warp-exponent argument is preferred.

% ============================================================
% SYMBOLIC VERIFICATION via src/sugra  (SymPy)
%   scripts:  src/examples/R_reduction_check.py
%             src/examples/F1_12d_uplift_check.py
% ============================================================

\section*{Symbolic verification (Sec.~6.2--6.3) with \texttt{src/sugra}}

\subsection*{(I) Does the 12d repackaged action reduce to the 10d IIB action?}

\textbf{Yes.} The non-trivial sector is gravity. For the block-diagonal ansatz
$\hat g = g_{\rm base}\oplus\hat g_{ab}(\tau)$ with the unit-determinant coset block
$\hat g_{ab}=\tfrac{1}{\tau_2}\!\left(\begin{smallmatrix}1&\tau_1\\\tau_1&\tau_1^2+\tau_2^2\end{smallmatrix}\right)$,
the 12d Ricci scalar must satisfy
\begin{equation}
\hat R = R_{\rm base} - \tfrac12 \frac{|\partial\tau|^2}{\tau_2^2}.
\end{equation}
Checked symbolically (\texttt{R\_reduction\_check.py}) with a flat base ($R_{\rm base}=0$) and the
algebraically-independent profiles $\tau_1=\sin x_0$, $\tau_2=e^{x_0}$: \texttt{SymPy} returns
\[
\hat R = -\tfrac12\big(\cos^2 x_0 e^{-2x_0}+1\big),\qquad
\hat R - \Big(\!-\tfrac12\tfrac{(\partial\tau_1)^2+(\partial\tau_2)^2}{\tau_2^2}\Big) = 0 .
\]
The coefficient is exactly $\tfrac12$. The remaining sectors were checked by hand:
$|\hat F_4|^2 = e^{-\Phi}|H_3|^2+e^{\Phi}|F_3-CH_3|^2$,
$|\tilde{\hat F}_7|^2=|\tilde F_5|^2$, and
$\tfrac18\!\int\hat C_4\wedge\hat F_4\wedge\hat F_4 \to -\tfrac{1}{4\kappa^2}\!\int C_4\wedge H_3\wedge F_3$.
So the Sec.~6.2 action reduces to the Einstein-frame IIB action ($g_s=1$).

\subsection*{(II) Do the brane uplifts solve the 12d Einstein equation of that action?}

For F1, D1, NS5, D5 one has $\tilde{\hat F}_7=0$ and the CS term vanishes, so the Sec.~6.2
action is pure 12d gravity $+\hat F_4$, with Einstein equation $R_{MN}=S_{MN}[\hat F_4]$
($n=4$, $D_{\rm trace}=12$). Results (\texttt{F1\_12d\_uplift\_check.py}):

\begin{itemize}
\item[(B)] \textbf{Control.} The genuine 12d electric membrane (warps $H^{-2/3}$ on $(t,x_1,z_1)$,
$H^{2/7}$ on the $9$ transverse dims, flux $\tfrac{2}{\sqrt\Delta}H^{-1}$ with $\Delta=\tfrac{21}{5}$,
$H$ harmonic in $9$) satisfies $R_{MN}=S_{MN}$ on \emph{all $12$ components}. (Validates the code.)
\item[(A)] \textbf{F1 uplift} (warps $H^{-3/4}|H^{1/4}|H^{-1/2}|H^{1/2}$, flux $H^{-1}$, $H$ harmonic
in $8$): \emph{mismatch on every component}. E.g. $R_{tt}=\tfrac{3H'^2}{8H^3}$ vs
$S_{tt}=\tfrac{7H'^2}{20H^3}$; and the torus components $R_{z_1}=-\tfrac{H'^2}{4H^{11/4}}$,
$R_{z_2}=+\tfrac{H'^2}{4H^{7/4}}$ are not matched by $S_{z_az_b}$. \emph{The F1 uplift does not
solve the 12d Einstein equation.}
\item[(C)] \textbf{F1 uplift with 10d weights} ($n_{\rm eff}=3$, $D_{\rm trace}=10$): the
$t,x_1,y_i$ components now \emph{match}, but $z_1,z_2$ still fail. I.e. the uplift solves the
\emph{10d} Einstein equations; the only equations it violates are the torus-direction components
$\hat R_{z_az_b}=S_{z_az_b}$, which have no 10d counterpart.
\end{itemize}

\textbf{Conclusion.} The Sec.~6.2 action reduces to IIB \emph{as a 10d functional} (torus frozen),
and the uplifts solve those 10d equations. They do \emph{not} solve the 12d Einstein equation
obtained by varying $\hat g_{MN}$ as a dynamical 12d metric; the failure is localized precisely to
the torus directions. This is the sharp, verified form of the ``repackaging, not a true 12d theory''
statement.

\subsection*{(III) Decomposing the torus-block failure (\texttt{F1\_torus\_decomp.py})}

The torus metric depends on the moduli only through $\hat g_{ab}(\tau)$, so the torus-block Einstein
equations are not independent of the axio-dilaton EOM: $\delta S/\delta\tau_i =
(\delta S/\delta\hat g_{ab})(\partial\hat g_{ab}/\partial\tau_i)$. There are three torus components
($z_1z_1,z_1z_2,z_2z_2$) but only two moduli $\tau_1,\tau_2$ plus the volume $\sigma$
($\det\hat g_{ab}=e^{2\sigma}$) that the ansatz freezes to $1$. For the F1 uplift ($\tau_1=0$,
genuine $n=4$, $D=12$):
\[
E_{z_1z_1}=R_{z_1z_1}-T_{z_1z_1}=\frac{H'^2}{10 H^{11/4}},\qquad
E_{z_2z_2}=\frac{H'^2}{10 H^{7/4}},
\]
and, with $g^{z_1z_1}=H^{1/2}$, $g^{z_2z_2}=H^{-1/2}$,
\[
\underbrace{g^{z_1z_1}E_{z_1z_1}-g^{z_2z_2}E_{z_2z_2}}_{\text{volume-preserving }(\tau_2)\text{ direction}} = 0,
\qquad
\underbrace{g^{z_1z_1}E_{z_1z_1}+g^{z_2z_2}E_{z_2z_2}}_{\text{volume-changing (radion) direction}} = \frac{H'^2}{5 H^{9/4}}\neq 0.
\]
The volume-preserving combination is the 10d axion+dilaton EOM (the trace term cancels here, so it is
weight-independent), and it \emph{holds}. The single failing equation is the torus-volume (radion)
modulus, sourced by $|F|^2\sim H'^2$. Freezing $\det\hat g_{ab}=1$ is thus not a consistent
truncation, and that one scalar equation is the entire obstruction to the uplift being a genuine 12d
solution; the axio-dilaton sector itself is perfectly consistent.

The same structure holds with the 10d trace weights ($n_{\rm eff}=3$, $D_{\rm trace}=10$): there
$E_{z_az_a}=H'^2/(8H^{\cdots})$, the volume-preserving combination is again $0$, and the radion
combination is $H'^2/(4H^{9/4})$. Only the radion coefficient depends on the weights; the dilaton
combination vanishes either way.

\paragraph{Can the mismatch be reconciled?} Partly. The volume-preserving (axio-dilaton) part is
\emph{exactly} the satisfied 10d $\tau$-EOM, so that part is reconciled. The irreducible piece is the
single radion equation $\delta S/\delta\sigma\propto H'^2\neq 0$, which has no 10d counterpart. It is
reconciled \emph{iff} the torus volume is treated as non-dynamical (the Sec.~6.2 stance: the action is
a 10d functional, $\sigma$ frozen), in which case $\delta S/\delta\sigma=0$ is simply not imposed and
the uplift is a bona fide solution of the actual (10d) variational problem; the nonzero residual is
then just the constraint force on a frozen modulus. If instead $\sigma$ is dynamical, the radion EOM
must hold, it does not, and the $10+2$ ansatz is not a consistent truncation. This is the precise
boundary between ``a 10d solution repackaged in 12d notation'' and ``a genuine 12d solution.''

\subsection*{(IV) Discussion of the 10d-weight check (C)}

The diagnostic that makes the structure transparent compares the \emph{full 12d} Ricci tensor
$R_{MN}$ of the F1 uplift against the flux stress tensor evaluated with \emph{10d weights}
($n_{\rm eff}=3$, $D_{\rm trace}=10$, $\Phi=0$). The motivation for these weights is that
$\hat F_4 = H_3\wedge dz_1$ carries only the degrees of freedom of a 10d 3-form, so the trace term
$\tfrac{n_{\rm eff}-1}{D_{\rm trace}-2}|F|^2$ natural to it is the 10d-3-form one. The component
output is
\begin{verbatim}
  [OK] t    R = 3*H'^2/(8*H^3)
  [OK] x1   R = -3*H'^2/(8*H^3)
  [OK] y0..y7  R = (H'^2 r^2 - 4 H'^2 y_i^2)/(8 H^2 r^2)
  [XX] z1   R = -H'^2/(4*H^(11/4))
  [XX] z2   R = +H'^2/(4*H^(7/4))
\end{verbatim}

\paragraph{The ten physical directions pass.} For $t,x_1,y_0,\dots,y_7$ one finds $R_{MN}=T_{MN}$
exactly. These are the ten directions that exist in 10d, and the equation they satisfy is the
\emph{10d type IIB Einstein equation} for the F1 string: the dilaton stress
$\tfrac12\partial_m\Phi\partial_n\Phi$ that would appear on the matter side in 10d is here supplied
by the curvature of the (geometrized) torus inside $R_{MN}^{(12)}$, while the flux supplies the
$H_3$ stress. So this block is a direct check that the uplift's 10d content is precisely the
ordinary 10d F1 solution.

\paragraph{What the 10d weights are telling us.} Imposing the \emph{genuine} 12d form-field
Einstein equation ($n=4$, $D_{\rm trace}=12$) component by component is the \emph{wrong} variational
problem for a constrained embedding: it treats $\hat g_{MN}$ as twelve independent fields, which the
ansatz does not (the torus block is the fixed function $\hat g_{ab}(\tau)$, $\det=1$). Against that
criterion (run (A)) the uplift fails on all twelve components. The correct procedure is to vary the
\emph{embedded} fields $g_{mn},\tau,B_2,\dots$; doing so returns the 10d IIB equations of motion,
which the brane solutions satisfy (verified directly in 10d, \texttt{full\_recheck.py}). The
``$10$d weights'' that make ten components agree in the diagnostic (C) are therefore not an ad hoc
concession but a reflection of the same fact: $\hat F_4=H_3\wedge dz_1$ contracts through its single
torus leg and gravitates as a 10d 3-form (cf.\ the genuine 12d membrane's inequivalent warps
$-2/3, 2/7$). With the embedded variation the ten physical equations and the axio-dilaton EOM hold;
only the radion --- a direction the embedding does not contain --- is left over.

\paragraph{The two torus directions fail --- but reducibly.} $z_1,z_2$ have no 10d counterpart, so
$R_{z_az_b}=T_{z_az_b}$ are genuinely new equations. As shown in (III) they decompose as
\[
g^{z_1z_1}E_{z_1z_1}-g^{z_2z_2}E_{z_2z_2}=0\ \ (\text{axio-dilaton EOM, satisfied}),\qquad
g^{z_1z_1}E_{z_1z_1}+g^{z_2z_2}E_{z_2z_2}=\tfrac{H'^2}{4H^{9/4}}\neq 0\ \ (\text{radion}).
\]
The opposite signs of $R_{z_1}<0$ and $R_{z_2}>0$ are the imprint of the two radii
$g_{z_1z_1}=H^{-1/2}, g_{z_2z_2}=H^{1/2}$ scaling inversely (unit volume), i.e. of the geometrized
dilaton. The mismatch is therefore \emph{not} a sickness of the solution and \emph{not} a failure of
the axio-dilaton sector: it is entirely the equation of motion of the frozen torus-volume modulus.

\paragraph{How it is reconciled.} The radion residual is $\delta S/\delta\sigma \propto H'^2$ --- the
``force'' on the volume modulus $\sigma$, switched on by the brane gradient. Whether this is a true
inconsistency depends on a single modelling choice:
\begin{itemize}
\item[(i)] \textbf{Torus volume non-dynamical} (the Sec.~6.2 stance: after the $\int_{T_2}$ the action
is a 10d functional and $\sigma$ is fixed by $\det\hat g_{ab}=1$). Then $\delta S/\delta\sigma=0$ is
\emph{never imposed} --- it is not one of the equations of the variational problem --- so there is no
inconsistency at all, and the uplift is a bona fide solution of the actual (10d) equations. The
nonzero residual is then merely the constraint force holding the frozen modulus in place. \emph{This
is how the mismatch is reconciled.}
\item[(ii)] \textbf{Torus volume dynamical} (treat Sec.~6.2 as genuine 12d gravity). Then
$\delta S/\delta\sigma=0$ must hold; it does not; so the uplift is not a solution and the $10+2$
ansatz is not a consistent truncation --- $\sigma$ would roll, sourced by $H'^2$.
\end{itemize}
As $H'\to 0$ (no brane) the radion residual vanishes, so the obstruction is sourced entirely by the
brane: the flat-torus vacuum is fine, and only the brane gradient excites the frozen modulus.

\paragraph{Summary: what the quotes mean.} The upshot is constructive. Substituting the $10+2$
ansatz into the Sec.~6.2 action returns the IIB action (sector by sector, (I)), so varying the
embedded fields $g_{mn},\tau,B_2,C_2,C_4$ gives the IIB equations of motion --- which every IIB brane
solution satisfies. \emph{This is precisely the content of the quotes in ``12d uplift'' and ``12d
solution'':} the configurations solve the equations of the embedded (constrained) action, equivalently
the 10d IIB EOM, but not the equations of \emph{unconstrained} 12d gravity. The latter contain one
extra equation with no 10d counterpart --- the radion $\delta S/\delta\sigma\propto|\hat F_4|^2$,
conjugate to the torus volume the embedding holds fixed. The recipe is uniform, and the radion source
$|\hat F_4|^2$ sorts its images into two tiers:
\begin{itemize}
\item \textbf{No 4-form flux} (axio-dilaton sector: D7, D(-1)) --- the radion source vanishes, so the
uplift solves even the unconstrained 12d equations: a \emph{genuine} 12d solution (the KK-monopole and
pp-wave of Sec.~4).
\item \textbf{With 4-form flux} (F1, D1, NS5, D5) --- the radion is sourced, so the uplift is an
effective ``12d solution'' in the embedded sense only.
\end{itemize}
Thus the quotes are doing exactly the work established above: ``$12$d solution'' $=$ solution of the
embedded action ($=$ IIB solution), promoted to a genuine 12d solution precisely when no 4-form flux
is excited.

% ============================================================
% VERIFICATION (sympy, anaconda python) of the 5-point obstruction (Sec. Effective Actions)
% Metric: ds^2 = eta_{mn}dx^m dx^n + g_ab(tau) dz^a dz^b,
%   g_ab = (1/tau2)[[1, tau1],[tau1, tau1^2+tau2^2]], tau1,tau2 functions of x^m.
%
% (1) The 12d Riemann components R_{ambn} were recomputed from the metric and AGREE
%     with the table in the main text, e.g.:
%       R_{1m1n}: dd2 -> 1/(2 t2^2),  (d t1)^2 -> 1/(4 t2^3),  (d t2)^2 -> -3/(4 t2^3)
%       R_{2m2n}: dd1 -> -t1/t2, dd2 -> (t1^2-t2^2)/(2 t2^2),
%                 (d t1)^2 -> (t1^2-3 t2^2)/(4 t2^3), (d t2)^2 -> (t2^2-3 t1^2)/(4 t2^3)
%     (all four rows checked; full agreement).
%
% (2) No-go check. Most general 2-derivative tau-scalar from one Riemann:
%       S = c^{ab} g^{mn} R_{ambn},  c^{ab} symmetric (alpha,beta,gamma).
%     Imposing that the box-tau1 and box-tau2 pieces vanish (1 constraint each) leaves a
%     1-parameter family for which sympy gives:
%       coeff (d t1)^2 = coeff (d t2)^2 = beta/(2 t1 t2),  coeff (d t1.d t2) = 0
%     i.e. S \propto (d tau1)^2 + (d tau2)^2 (SAME sign, no cross term) = U(1)-invariant |d tau|^2.
%     Hence the charge-2 combination P_m P^m \propto (d t1)^2 - (d t2)^2 + 2i d t1 . d t2
%     (opposite signs + cross term) is NOT representable. Confirms the obstruction.
% (scripts: /tmp/riem_check.py, /tmp/obstruction_check2.py)
% ============================================================

% ============================================================
% DETAILED KK-loop R^4 derivation (for Section 7.2, to be truncated)
% Goal: extract the R^4 coefficient from the existing 4-point KK-loop
%       integral \eqref{4-point Schwinger parameter}, following the
%       standard procedure (low-energy expansion -> simplex -> proper-time
%       integral -> lattice sum -> Eisenstein functional equation),
%       and identify the dimension dependence.
% ============================================================

\section*{Detailed KK-loop $R^4$ coefficient (working)}

\paragraph{Starting point.}
The existing 4-point KK-loop integral (Section 7.2) is
\begin{equation}
\mathcal{A}_4 \propto \sum_{m,n}\int_{\Lambda^{-2}}^\infty \frac{d\lambda}{\lambda^{3/2}}\,
e^{-\lambda M_{m,n}^2}\, P(s,t;\lambda),\qquad
M_{m,n}^2 = \frac{|n-m\tau|^2}{R^2\tau_2^2},
\end{equation}
\begin{equation}
P(s,t;\lambda) = \int_0^1 d\rho_3\int_0^{\rho_3}d\rho_2\int_0^{\rho_2}d\rho_1\,
e^{-\lambda M(s,t;\rho)},\qquad
M(s,t;\rho) = s\rho_1\rho_2 + t\rho_2\rho_3 + u\rho_1\rho_3 + t(\rho_1-\rho_2).
\end{equation}

\paragraph{Step 1: low-energy expansion (this gives the local terms).}
The local higher-derivative terms come from expanding $P$ in powers of the external momenta:
\begin{equation}
P(s,t;\lambda) = \int_\Delta d^3\rho\, e^{-\lambda M(s,t;\rho)}
= \sum_{k=0}^\infty \frac{(-\lambda)^k}{k!}\int_\Delta d^3\rho\, M(s,t;\rho)^k,
\end{equation}
where $\Delta = \{0\le\rho_1\le\rho_2\le\rho_3\le1\}$.
Since $M \sim (\text{Mandelstam}) \sim (\text{momentum})^2$, the order-$k$ term carries $2k$ extra derivatives and produces the $D^{2k}R^4$ contact term.
The leading ($k=0$) term is the $R^4$ coefficient, equal to the simplex volume
\begin{equation}
P(0,0;\lambda) = \int_\Delta d^3\rho
= \int_0^1 d\rho_3\,\frac{\rho_3^2}{2} = \frac{1}{3!} = \frac{1}{6}.
\end{equation}

\paragraph{Step 2: the $R^4$ coefficient as a regulated proper-time sum.}
Hence
\begin{equation}
\mathcal{C} \;\propto\; \sum_{(m,n)\neq(0,0)} \int_{\Lambda^{-2}}^\infty \frac{d\lambda}{\lambda^{3/2}}\, e^{-\lambda M_{m,n}^2}.
\end{equation}
The massless mode $(m,n)=(0,0)$ is omitted: with $M=0$ the integral is IR divergent and gives the
non-analytic supergravity exchange ($1/stu$), not a local term.
The $\lambda\to0$ end is UV divergent (regulated by $\Lambda$); that divergence is a local $R^4$
counterterm. The renormalised finite part is the analytic continuation of
\begin{equation}
\int_0^\infty d\lambda\, \lambda^{s-1} e^{-\lambda M^2} = \Gamma(s)\,(M^2)^{-s}.
\end{equation}

\paragraph{Step 3: the measure fixes $s$.}
For a $d$-dimensional loop the box proper-time measure is $\lambda^{3-d/2}$
($\lambda^3$ from the Jacobian of the four Feynman parameters $\lambda_i=\lambda\rho_i$, and
$\lambda^{-d/2}$ from the loop-momentum Gaussian $\int d^dk\,e^{-\lambda k^2}$).
The quoted measure $\lambda^{-3/2}$ is therefore $d=9$:
\begin{equation}
3-\tfrac{d}{2} = -\tfrac32 \;\Rightarrow\; d=9,\qquad
s \equiv 4-\tfrac{d}{2} = -\tfrac12,\ \text{consistent with}\ \lambda^{s-1}=\lambda^{-3/2}\Rightarrow s=-\tfrac12 .
\end{equation}
[NOTE/CHECK: the integrand power is $s-1 = -3/2$, i.e. $s=-1/2$, which is the value entering
$\int d\lambda\,\lambda^{s-1}e^{-\lambda M^2}=\Gamma(s)(M^2)^{-s}$. The ``$4-d/2$'' bookkeeping
($\nu=4$ propagators, $\int d^dp\,(p^2+M^2)^{-4}$) gives the same $s=-1/2$ at $d=9$. Tension to flag
to author: the setup writes a 10d loop momentum $\nabla^2_{10}$, which would be $d=10$,
$\lambda^{-2}$, $s=-1$; but the measure as written is $\lambda^{-3/2}$, i.e. $d=9$.]
With $s=-\tfrac12$,
\begin{equation}
\mathcal{C} \;\propto\; \Gamma\!\left(-\tfrac12\right)\sum_{(m,n)\neq(0,0)} (M_{m,n}^2)^{1/2}.
\end{equation}

\paragraph{Step 4: lattice sum $\to$ Eisenstein.}
With $(M_{m,n}^2)^{1/2} = |n-m\tau|/(R\tau_2)$ and the bijection $(m,n)\mapsto(n,-m)$ of
$\mathbb{Z}^2\setminus\{0\}$ (so $\sum|n-m\tau|^a = \sum|m+n\tau|^a$),
\begin{equation}
\sum_{(m,n)\neq(0,0)} (M_{m,n}^2)^{1/2}
= \frac{1}{R\tau_2}\sum_{(m,n)\neq(0,0)} |m+n\tau|
= \frac{1}{R}\,\tau_2^{-1/2}\, E_{-1/2}(\tau),
\end{equation}
where $E_s(\tau)\equiv\sum_{(m,n)\neq(0,0)}\tau_2^{s}/|m+n\tau|^{2s}$, used at $s=-\tfrac12$ so that
$|m+n\tau|^{-2s}=|m+n\tau|$ and the extracted $\tau_2^{-s}=\tau_2^{1/2}$ combines with the $1/\tau_2$
to give $\tau_2^{-1/2}$.

\paragraph{Step 5: functional equation $\to f_0=E_{3/2}$.}
The completed series $E_s^\ast(\tau)\equiv\pi^{-s}\Gamma(s)\zeta(2s)E_s(\tau)$ obeys $E_s^\ast=E_{1-s}^\ast$.
At $s=-\tfrac12$ ($1-s=\tfrac32$), with $\Gamma(-\tfrac12)=-2\sqrt\pi$, $\Gamma(\tfrac32)=\tfrac{\sqrt\pi}{2}$, $\zeta(-1)=-\tfrac{1}{12}$:
\begin{equation}
\underbrace{\pi^{1/2}(-2\sqrt\pi)(-\tfrac{1}{12})}_{=\pi/6}\,E_{-1/2}
= \underbrace{\pi^{-3/2}\tfrac{\sqrt\pi}{2}\zeta(3)}_{=\zeta(3)/2\pi}\,E_{3/2}
\;\Rightarrow\;
E_{-1/2} = \frac{3\zeta(3)}{\pi^2}\,E_{3/2} = \frac{3\zeta(3)}{\pi^2}\, f_0 .
\end{equation}

\paragraph{Result.}
Collecting, and writing $R^{-1}\tau_2^{-1/2}=(R^2\tau_2)^{-1/2}=A^{-1/2}$ with $A=R^2\tau_2$ the torus area,
\begin{equation}
\boxed{\;\mathcal{C} \;\propto\; \Gamma\!\left(-\tfrac12\right)\frac{3\zeta(3)}{\pi^2}\, A^{-1/2}\, f_0(\tau,\bar\tau)\;}
\end{equation}
The $\tau$-dependence is entirely in $f_0=E_{3/2}$; the prefactor $A^{-1/2}$ is independent of the
complex structure at fixed area, so the result is modular invariant.
This matches the zero-mode part of the perturbative effective action (Section 7.1):
$f_0 = 2\zeta(3)\tau_2^{3/2}+\tfrac{2\pi^2}{3}\tau_2^{-1/2}+(\text{D(-1) instantons})$.

\paragraph{Diagnostic ($d=9$ vs $d=10$).}
Only the $d=9$ measure ($s=-\tfrac12$) reflects onto $E_{3/2}=f_0$.
A genuine $d=10$ loop would carry $\lambda^{-2}$, i.e. $s=-1$, which $E_s^\ast=E_{1-s}^\ast$ maps to
$E_2$ (weight 2, not invariant), and $\Gamma(4-\tfrac d2)=\Gamma(-1)$ has a pole (the log running of
$R^4$ in even dimensions). [Cross-check with author's functional-determinant computation: the genuine
$d=10$ loop there ($\int d^{10}p$) gives the modular-COVARIANT $E_6$ for the vacuum energy -- a
different observable, but the same lesson: a true 12d/$d=10$ loop does not give the invariant $f_0$.]
% ============================================================
